\documentclass[tikz,border=10pt]{standalone}
\usepackage[UTF8]{ctex}
\usetikzlibrary{arrows.meta,backgrounds,fit,positioning,shadows.blur,shapes.geometric}

\definecolor{ink}{HTML}{1F2937}
\definecolor{muted}{HTML}{5B6472}
\definecolor{lineblue}{HTML}{2563EB}
\definecolor{frontendfill}{HTML}{EFF6FF}
\definecolor{backendfill}{HTML}{F0FDF4}
\definecolor{datafill}{HTML}{FFF7ED}
\definecolor{actorfill}{HTML}{F8FAFC}
\definecolor{groupborder}{HTML}{CBD5E1}

\tikzset{
  font=\sffamily,
  >={Latex[length=3mm,width=2mm]},
  actor/.style={draw=ink, rounded corners=2pt, fill=actorfill,
    minimum width=28mm, minimum height=13mm, align=center,
    line width=.8pt, font=\bfseries\sffamily},
  component/.style={draw=ink, rounded corners=2pt, fill=white,
    minimum width=40mm, minimum height=18mm, text width=36mm,
    align=center, inner sep=4pt, line width=.8pt,
    font=\small\sffamily, blur shadow={shadow opacity=.10}},
  datastore/.style={cylinder, shape border rotate=90, aspect=.27,
    draw=ink, fill=white, minimum width=40mm, minimum height=19mm,
    align=center, line width=.8pt, font=\small\sffamily},
  dependency/.style={-Latex, draw=muted, line width=.8pt},
  request/.style={-Latex, draw=lineblue, line width=1.1pt},
  labeltext/.style={font=\scriptsize\sffamily, fill=white,
    inner sep=1.5pt, text=muted},
  group/.style={draw=groupborder, rounded corners=5pt,
    line width=.8pt, inner sep=2mm},
  grouptitle/.style={font=\bfseries\small\sffamily, text=ink,
    fill=white, inner sep=2pt},
  note/.style={draw=groupborder, rounded corners=2pt, fill=actorfill,
    text width=325mm, align=left, inner sep=5pt,
    font=\small\sffamily, text=ink}
}

\newcommand{\componentglyph}[1]{%
  \begin{scope}[shift={(#1.north east)},xshift=-2.5mm,yshift=-1.5mm,
      draw=ink,line width=.55pt]
    \draw (-3mm,-1.7mm) rectangle (0,1.7mm);
    \draw (-4.2mm,0.6mm) rectangle (-2.5mm,1.55mm);
    \draw (-4.2mm,-1.55mm) rectangle (-2.5mm,-0.6mm);
  \end{scope}%
}

\newcommand{\diagramtitle}[2]{%
  \node[font=\bfseries\LARGE\sffamily, text=ink] (title) at (15.5,9.5) {#1};
  \node[below=1.5mm of title, font=\small\sffamily, text=muted] {#2};
}

\newcommand{\groupcaption}[2]{%
  \node[grouptitle, anchor=south west]
    at ([xshift=2mm,yshift=1mm]#1.north west) {#2};
}

\begin{document}
\begin{tikzpicture}
  \diagramtitle{UC08 直播间管理——UML 组件图}
    {展示直播间创建与状态维护、OBS 推流、SRS 分发、观众播放和直播评论所依赖的组件}

  \node[actor] (anchor) at (0,5.7) {作者／主播};
  \node[actor] (viewer) at (0,1.8) {观众};
  \node[component] (studio) at (4,5.7) {\textbf{StudioLiveView}\\\footnotesize 创建、开播与停播};
  \node[component] (liveView) at (4,1.8) {\textbf{LiveListView / LiveRoomView}\\\footnotesize 列表、播放与评论};
  \node[component] (client) at (8.5,3.75) {\textbf{live.js + livePlayer.js}\\\footnotesize 房间 API 与 FLV/HLS 播放};

  \node[component] (security) at (13,3.75) {\textbf{SecurityConfig + JWT Filter}\\\footnotesize 公开观看／主播操作校验};
  \node[component] (controller) at (17.5,5.2) {\textbf{LiveRoomController}\\\footnotesize 房间、状态与直播评论};
  \node[component] (service) at (22,5.2) {\textbf{LiveRoomService}\\\footnotesize 推流密钥与状态规则};
  \node[component] (task) at (19.75,1.5) {\textbf{LiveMaintenanceTask}\\\footnotesize 弹幕与过期房间清理};

  \node[component] (mapper) at (26.5,5.2) {\textbf{Mapper 层}\\\footnotesize LiveRoomMapper · CommentMapper};
  \node[datastore] (mysql) at (31,5.2) {\textbf{MySQL 8}\\\footnotesize live\_rooms / comments};
  \node[component] (obs) at (26.5,1.2) {\textbf{OBS}\\\footnotesize RTMP 推流客户端};
  \node[component] (srs) at (31,1.2) {\textbf{SRS 媒体服务器}\\\footnotesize RTMP 输入 · FLV/HLS 输出};

  \draw[request] (anchor) -- node[labeltext,above]{管理房间} (studio);
  \draw[request] (viewer) -- node[labeltext,above]{观看／评论} (liveView);
  \draw[request] (studio) -- (client);
  \draw[request] (liveView) -- (client);
  \draw[request] (client) -- node[labeltext,above]{HTTP + JWT} (security);
  \draw[request] (security) -- (controller);
  \draw[dependency] (controller) -- (service);
  \draw[dependency] (service) -- (mapper);
  \draw[dependency] (task) to[out=15,in=-135] (mapper);
  \draw[dependency] (mapper) -- (mysql);
  \draw[request] (anchor.south) |- node[labeltext,pos=.72,below]{采集音视频} (obs.south);
  \draw[request] (obs) -- node[labeltext,above]{RTMP} (srs);
  \draw[request] (srs.west) to[out=160,in=-20] node[labeltext,below]{FLV / HLS} (liveView.east);

  \foreach \n in {studio,liveView,client,security,controller,service,task,mapper,obs,srs}{\componentglyph{\n}}
  \begin{scope}[on background layer]
    \node[group, fill=frontendfill, fit=(studio)(liveView)(client)] (fg) {};
    \node[group, fill=backendfill, fit=(security)(controller)(service)(task)] (bg) {};
    \node[group, fill=datafill, fit=(mapper)(mysql)(obs)(srs)] (dg) {};
  \end{scope}
  \groupcaption{fg}{Vue 3 前端}
  \groupcaption{bg}{Spring Boot 后端}
  \groupcaption{dg}{数据与直播媒体基础设施}
  \node[note, anchor=north west] at (-1.4,-0.8) {\textbf{职责边界：}\quad Spring Boot 只维护房间、开播状态和评论；OBS 负责采集并通过 RTMP 推流；SRS 负责媒体接收与 FLV/HLS 分发。};
\end{tikzpicture}
\end{document}
